\section{Semantic Annotation}
\label{app:annotation}

\paragraph{Label sets}
The annotator outputs a JSON object. \spin{} uses three of its fields:
\begin{itemize}
  \setlength\itemsep{0pt}
  \item \textbf{Reply relation} (11 values): \texttt{agree\_with\_parent}, \texttt{disagree\_with\_parent}, \texttt{mock\_parent}, \texttt{sarcastic\_agreement}, \texttt{clarify\_parent}, \texttt{supplement\_parent}, \texttt{question\_parent}, \texttt{answer\_parent}, \texttt{topic\_shift}, \texttt{direct\_comment\_to\_target}, \texttt{unknown}.
  \item \textbf{Propagation type} (9 values): \texttt{stance\_inheritance}, \texttt{stance\_reversal}, \texttt{sarcastic\_reversal}, \texttt{stance\_amplification}, \texttt{stance\_softening}, \texttt{direct\_stance\_expression}, \texttt{target\_shift}, \texttt{unclear\_propagation}, \texttt{no\_parent\_context}.
  \item \textbf{Target relation} (10 values): \texttt{same\_target}, \texttt{attribute\_of\_target}, \texttt{consequence\_of\_target}, \texttt{reason\_for\_target}, \texttt{opposition\_to\_related\_target}, \texttt{support\_to\_related\_target}, \texttt{reply\_author\_only}, \texttt{topic\_shift}, \texttt{irrelevant}, \texttt{unclear}.
\end{itemize}
The remaining fields (e.g., an auxiliary stance guess and a confidence score) are not used by \spin{}. Out-of-vocabulary values are mapped to the default of each field.

\paragraph{Prompt}
Appendix~\ref{app:prompt} gives the system prompt and the user prompt template, translated from the original Chinese; the Chinese originals are released with our code. Slots in braces are filled per node: the target anchor is truncated to 300 characters and the causal context to 600 characters. The template also contains slots for downstream replies and for a post-hoc context; in the causal mode used throughout this paper these slots are replaced by the fixed strings shown in Appendix~\ref{app:prompt}, so the annotator never sees any descendant of the node. Requests are sent with temperature 0.1 and a maximum of 800 output tokens; a response that is not valid JSON is retried up to three times.

\paragraph{Agreement with gold relations}
On edges whose parent holds a gold stance and whose propagation type maps onto a relation (inheritance and amplification $\rightarrow$ inherit; reversal and sarcastic reversal $\rightarrow$ reverse; softening $\rightarrow$ neutralize), the annotation agrees with the gold relation for about 56\% of edges on C-MTCSD and 52\% on ZS-CSD.

\section{Data Processing}
\label{app:data}

\paragraph{C-MTCSD}
We build one tree per conversation from the official \texttt{train/valid/test} files and keep the official split of every post and reply. Two steps are needed to cover the full split. First, the release stores root posts as conversation context rather than as nodes; we add the 452 root posts as depth-0 nodes, so that depth-1 comments have a parent in the tree. Root posts receive the propagation type \texttt{no\_parent\_context} by rule, without an LLM call, and are therefore affected only by absolute knowledge during inference. Second, we make sure that every reply is annotated: 219 replies were missing from an earlier run of our annotation pipeline (identifiers repeat across topics, so a resumed run skipped some of them, and a few requests were rejected by the API), and we re-annotated them. All but one now carry an annotation; the remaining reply is still rejected by the content filter and is kept with an \emph{unknown} key, which the backoff chain handles. The resulting splits contain 17,102 / 3,663 / 3,498 nodes, identical to the official ones, so our test scores are directly comparable with published results.

\paragraph{ZS-CSD}
The release stores one instance per root-to-node path. We reconstruct trees by matching each path's ancestor texts, assign node identifiers by position, and keep the original split and target of every node. Root posts, which have no parent, receive the propagation type \texttt{no\_parent\_context}.

\section{Hyperparameters and Robustness}
\label{app:hyper}

\paragraph{Selected exponents}
Table~\ref{tab:lambda} lists the $(\lambda_r,\lambda_a)$ selected on validation. The selected values differ between models: the GPT models rely more on relative knowledge, while on C-MTCSD the Llama models rely more on absolute knowledge, which mainly corrects their over-prediction of \nonelabel{}. On ZS-CSD, where trees are deeper and parent stances matter more, larger values of $\lambda_r$ are selected. For GPT-4o on ZS-CSD (with both kNN and CoT prompting), $\lambda_a = 0$ is selected, so \spin{} reduces to relative knowledge only; for GPT-4o-mini full-path on C-MTCSD, $\lambda_r = 0$ is selected, so only absolute knowledge is used. SITPCL selects large exponents for all three seeds ($\lambda_r \geq 2$, $\lambda_a \geq 1$), yet the gains seen on validation do not carry over to test (Appendix~\ref{app:supervised}).

\paragraph{Grid}
The grid was first $\{0,\ldots,2\}$. Because GPT-4o-mini on ZS-CSD selected the boundary value $\lambda_r = 2$ on validation, we extended it to the grid above and used the extended grid for all settings. On C-MTCSD the selections for all four LLMs are unchanged; on ZS-CSD GPT-4o-mini moves from 46.01 (+2.30, 95\% CI [+0.41, +4.22]) to 46.91. For Llama-3.1-70B on ZS-CSD the validation surface is flat, and the selected (4, 1.5) is worse on test than relative knowledge alone (49.51).

\paragraph{Cross-validated selection}
As a robustness check we also select $(\lambda_r,\lambda_a)$ by tree-grouped 5-fold cross-validation within the validation split: for each fold, $E$ is estimated on the other four folds and inference is run on the held-out fold, and $\lambda$ is chosen by the pooled score. This protocol was proposed after observing the SITPCL results; before running it, we fixed the rule that it would replace the main protocol only if it did not lower any LLM setting and restored a gain for SITPCL. It did not restore a gain for SITPCL, so we keep the simpler protocol of Section~\ref{sec:hyper} in the main paper. Table~\ref{tab:cv} shows that all eight LLM settings remain significant under cross-validation, and four of them select exactly the same $\lambda$; GLAN reaches 65.83 topic-averaged \favg{} (64.04 without \spin{}).

\section{Supervised Models}
\label{app:supervised}

\paragraph{GLAN port}
Five changes were needed to run GLAN on C-MTCSD: model selection on the validation split only (the released code selected predictions by test score), left truncation, keeping the last 8 nodes of the path, Chinese targets with \texttt{chinese-roberta-wwm-ext}, and removing the experiment-tracking dependency.

\paragraph{Transfer diagnosis}
For each seed we compare three settings on the test split: the protocol ($E$ and $\lambda$ from validation), a $\lambda$ upper bound ($E$ from validation, $\lambda$ chosen on test) and a full upper bound ($E$ and $\lambda$ from test). For SITPCL the full upper bound is +2.88, +4.65 and +4.11 points while the protocol reaches $-0.32$, +1.55 and +1.97 (Fig.~\ref{fig:sitpcl}); the Jensen--Shannon divergence between the validation and test rows of $E$ is about 0.005, against about 0.001 for GPT-4o-mini.

\section{Results on the English MT-CSD}
\label{app:mtcsd}

MT-CSD~\cite{mtcsd} is an English multi-turn conversational stance dataset on five targets (Biden, Bitcoin, SpaceX, Tesla and Trump; 10,775 / 2,176 / 2,336 nodes for train / validation / test). We apply \spin{} with the same pipeline as on the Chinese datasets: the same semantic taxonomy and annotation prompt in the causal mode, a knowledge base built from the MT-CSD training split, GPT-4o-mini with kNN few-shot prompting as the base model, and $\lambda_r,\lambda_a$ selected on the validation split with the same grid.

Table~\ref{tab:mtcsd} shows the results. \spin{} improves pooled Macro-F1 by +1.48 points (95\% CI [+0.56, +2.38]; significant under all three tests), about a third of its gain on C-MTCSD with the same base model (+4.25). \favg{} does not improve significantly (+0.78, [$-0.21$, +1.74]), and topic-averaged \favg{} falls slightly ($-0.79$). The composition of the gain differs from the Chinese datasets. Relative knowledge alone brings nothing: validation selects $\lambda_r = 0$. The gain comes mostly from absolute knowledge (+1.17 alone), and combining both adds a little more. The semantic annotation still carries information: a key with the propagation type alone gives 61.44, whereas depth alone (59.82) and a random key (59.91) give no gain (we did not test these differences for significance). Below we discuss why the gain is smaller.

\paragraph{Semantic structure differs across languages}
The taxonomy of reply relations, propagation types and target relations, and the prompt that elicits them, were designed and refined on Chinese Weibo conversations. Chinese and English replies express stance toward a parent in different ways. Chinese Weibo replies are often very short and rely on rhetorical questions, irony and conventional expressions whose meaning depends on context, which is what the categories \emph{sarcastic reversal} and \emph{softening} were built to capture; English replies more often state disagreement explicitly or quote and paraphrase the parent. The cross-corpus summary is reported in the main paper; the following analysis discusses where the taxonomy stops fitting. Reversals are rare in MT-CSD: only 16.6\% of replies to a parent with a stance take the opposite stance, against 35.6\% in C-MTCSD, and even replies annotated as a \emph{reversal} keep their parent's stance more often than they reverse it (51\% versus 24\%). More tellingly, the irony category does not carry over. In C-MTCSD a \emph{sarcastic reversal} reverses the parent's stance in 40\% of cases, above the corpus base rate; in MT-CSD it does so in only 11\%, below the base rate, and instead points to a neutral reply (36\%). After removing the base rates, the log2-lift correlation between C-MTCSD and MT-CSD is 0.54 over all cells (0.80 when restricted to propagation types with at least 200 edges in both corpora), against 0.85 between the two Chinese corpora; the agreement is weakest for the reversal column (0.45, against 0.89 for inheritance and 0.82 for neutralization). Inheritance and neutralization thus transfer across languages, but the categories built to capture Chinese irony and rebuttal do not. When the categories fit the data less well, annotations become less consistent, keys group replies that propagate stance differently, and the retrieved distributions carry less information. Our analysis of the two Chinese corpora already shows that the base rates of the relations are corpus-specific (Section~\ref{sec:ana-gains}); across languages, the mapping from semantic description to propagation behaviour itself shifts.


\paragraph{Little correction room}
93.1\% of MT-CSD replies have their parent in the data, so the tree structure is available. Yet replacing the inferred parent posterior with the gold parent stance (the upper bound of Section~\ref{sec:ana-room}) gives 61.26, no higher than \spin{} itself (61.39), as on C-MTCSD. Because reversals are rare in MT-CSD, most replies simply share their parent's stance, and a base model that reads the parent in its context window already captures this, so knowing the parent's stance adds little. What remains for \spin{} to correct is mainly the base model's errors on \nonelabel{}, which absolute knowledge addresses.

\paragraph{Implication}
\spin{} still improves the English base model, and does not depend on the language in its inference; but its knowledge base inherits the language and platform of the taxonomy used to annotate it. Adapting the taxonomy and the annotation prompt to English conversational conventions, or learning the taxonomy from data, is a natural next step.

\section{Negative Results}
\label{app:negative}

The following variants did not yield stable gains and are not part of \spin{}. We report them so that others do not repeat them.

\paragraph{Cross-topic knowledge base}
On C-MTCSD we built the knowledge base for each topic from the training data of the other four topics. Validation selected $\lambda_r = \lambda_a = 0$, i.e., no correction, and forcing positive exponents lowered test scores. At the retrieval level the cross-topic knowledge base still predicts the relation better than the majority class (0.551 versus 0.449 with the propagation-type and depth key), so propagation patterns are partly shared across topics, but in inference this is outweighed by topic-specific base rates and topic-specific LLM biases. Calibrating the channel to the target topic did not help either: EM label-shift adjustment~\cite{saerens2002adjusting}, base-rate calibration from $k$ labelled target examples ($k$ from 10 to all, five seeds) and the training label distribution as prior all gave no gain. Only an oracle calibration with the test label distribution helped, which uses test information and only shows that the problem is the mismatch of distributions. These results are why \spin{} builds its knowledge base from in-domain training data.

\paragraph{Random keys}
Replacing the retrieval key by a random one makes validation select $\lambda$ close to zero, and forcing $\lambda = 2$ steadily lowers Macro-F1 (to 54.62 on an earlier version of our C-MTCSD data). The gains therefore come from the retrieved knowledge, not from the inference procedure itself.

\section{Additional Results}
\label{app:results}

\paragraph{Per-topic results on C-MTCSD}
Table~\ref{tab:topic} breaks down the gain for GPT-4o-mini kNN. \spin{} improves four of the five topics and leaves Non-marriage Doctrine unchanged; the largest gain is on Apollo Go, which is also the topic with the largest gain for every other LLM. For GLAN the largest gain is on iPhone 15 (Macro-F1 +4.3, +14.6 and +9.7 for the three seeds, all significant), and Apollo Go, where GLAN is already strong, barely changes. The \favg{} drop of Llama-3.1-70B (75.24 $\rightarrow$ 74.61 topic-averaged) comes from iPhone 15 and Non-marriage Doctrine.

\paragraph{Breakdowns on ZS-CSD}
For GPT-4o-mini kNN, the gain is larger for noun-phrase targets (+4.67 Macro-F1, 95\% CI [+0.99, +8.22]) than for targets phrased as a claim (+2.37, [$-0.16$, +5.12]). It is spread over depths without a clear trend (+4.27 at depth 2, $-0.21$ at depth 4, +4.08 at depth 5; no single depth is significant on its own). Nodes whose speaker has not appeared earlier in the path gain +3.81 ([+1.19, +6.37]), while replies by a speaker who already took part gain only +0.68 (not significant). Across the 45 test targets, 24 improve and 18 get worse. The base model often predicts \against{} for \favor{} nodes (396 cases); \spin{} raises \favor{} F1 from 0.343 to 0.434. For Llama-3.1-8B, 73\% of base predictions are \against{}, and \nonelabel{} F1 rises from 0.197 to 0.571.

\paragraph{Gains by depth}
The main paper reports the depth-wise gain breakdown; the following analysis discusses the observed pattern. On C-MTCSD \spin{} helps at every depth, and the gain grows with depth, from +3.5 points at depths 1--2 to +5.2--5.8 at depths 3--5, where the base model is also weakest. With CoT prompting on ZS-CSD the trend is clearer still: there is no gain at depth 2, and +5.4 to +7.6 points from depth 4 onward. For GPT-4o with CoT, \spin{} lowers Macro-F1 at depths 2 and 3 and raises it from depth 4 onward, so its small overall gain (Section~\ref{sec:ana-room}) hides opposite effects at shallow and deep levels. Only ZS-CSD with kNN prompting shows no trend. Deeper replies depend more on the conversation above them, which is where propagation knowledge adds the most. The per-depth samples are small, however, and we do not test the differences between depths for significance.

\section{Fine-tuned Base Models and Cost}
\label{app:finetune}

\paragraph{Fine-tuned Qwen2.5-7B with less training data}
To test whether \spin{} helps a fine-tuned model that has seen less data, we fine-tune Qwen2.5-7B on random samples of 10\%, 25\% and 50\% of the ZS-CSD training split (same recipe as in Section~\ref{sec:setup}; model selection on the full validation split). For a fair comparison, the knowledge base is built only from the training examples that the model was fine-tuned on. Table~\ref{tab:finetune} and Fig.~\ref{fig:finetune} show the results. \spin{} brings no gain at any size: validation selects relative knowledge only ($\lambda_r$ of 1--1.5), which raises the validation score by 0.4--1.5 points but lowers the test score by 0.6--1.3 (significant at 50\%). The gold-parent bound lies 3.7--4.6 points above the base for these models, so room exists, yet \spin{} cannot use it. Three explanations can be ruled out for the 10\% model. A larger knowledge base does not help: with the knowledge base built from the full training split, \spin{} still gives 53.00 against 53.35. The channel transfers well: $E$ is nearly identical on validation and test (row-wise Jensen--Shannon divergence at most 0.002), and even exponents chosen on test give no gain. Correlated errors are not the cause either: the risk ratio of child given parent errors is 1.16, close to that of GPT-4o-mini (1.12).

What distinguishes the fine-tuned models is how well their own predictions already capture propagation. Given the gold parent stance, the relation implied by a model's prediction (inherit, reverse or neutralize) is correct for 51.3\%, 56.5\%, 58.0\% and 61.3\% of replies with 10\%, 25\%, 50\% and 100\% of the training data, against 55.0\% for the relations retrieved from the knowledge base and only 41.7\% and 34.8\% for GPT-4o-mini and Llama-3.1-8B with kNN prompting (Fig.~\ref{fig:finetune}b). Absolute knowledge has nothing to correct either, since the fine-tuned models do not over-predict \nonelabel{}. On C-MTCSD, GPT-4o-mini's implied relations are already more accurate than the retrieved ones (68.4\% against 58.7\%; Section~\ref{sec:ana-prompt}), and the gain there comes mostly from correcting its over-prediction of \nonelabel{} (Section~\ref{sec:ana-errors}). In both cases \spin{} supplies statistics of the training data that the base model lacks; a model fine-tuned on that data already has them.

\paragraph{Cost}
Table~\ref{tab:cost} summarizes the cost of the annotation step. Annotating all splits of all three datasets takes about 56{,}000 GPT-4o-mini calls, about 1{,}200 input tokens each, and under US\$20 at list prices; the annotation is shared by all base models, and inference over a test split runs in under a second on a CPU. For comparison, fine-tuning Qwen2.5-7B with LoRA on the ZS-CSD training split takes about 1.1 hours on one A100 (40\,GB), and StanceAgent-R1 requires GPT-4o teacher trajectories followed by supervised and reinforcement-learning stages on two A800 GPUs~\cite{li2026stanceagent}.

\section{More Cases}
\label{app:cases}

Table~\ref{tab:more-cases} gives further cases: two corrections and two cascades of errors.

\paragraph{Corrections}
In C4 (topic Apollo Go, target: the Apollo Go robotaxi service), the reply answers a doubt about safety with a counter-question. The base model sees no explicit stance. The parent's stance is itself inferred wrongly (\nonelabel{}), so relative knowledge has little to pass on, and the correction must come mainly from absolute knowledge: replies annotated as a sarcastic reversal that mocks the parent's author are rarely neutral. In Z4 (target: \emph{parents should check students' homework}), the reply mirrors the parent sentence by sentence with the roles of teacher and parent swapped. The base model is misled by the shared negative wording and predicts the parent's stance; the annotated reversal of a correctly inferred parent gives the right answer.

\paragraph{Cascades}
C5 and Z5 show how one wrong inference spreads down a path. In C5 (target: quitting a job without a new one lined up), the base model correctly predicts \nonelabel{} for both replies, but the root post's stance is inferred wrongly, the first reply is annotated as a reversal and changed to \against{}, and this error is then passed on to its own reply. In Z5 (target: \emph{disliking babies with the XYY ``super-male'' genotype}), the parent of Z5a is neutral but inferred as \against{}; both Z5a and Z5b are annotated as reversals, so the error flips polarity twice, overturning two correct base predictions. As in the main text, the cascades occur on deeper paths of ZS-CSD and show that \spin{} is only as good as the inferred parent stance.

\section{Supplementary Figures}
\label{app:figures}

\begin{figure}[H]
\centering
\includegraphics[width=\columnwidth]{figures/fig_data_depth.pdf}
\caption{Depth distribution of the two datasets (root = 0).}
\label{fig:depth-dist}
\end{figure}

\begin{figure}[H]
\centering
\includegraphics[width=\columnwidth]{figures/fig_supervised_seeds.pdf}
\caption{Per-seed gain of \spin{} on the supervised models (pooled Macro-F1).}
\label{fig:seeds}
\end{figure}

\begin{figure}[H]
\centering
\includegraphics[width=\columnwidth]{figures/fig_sitpcl_transfer.pdf}
\caption{Transfer diagnosis for SITPCL: gain on the validation split where $\lambda$ is chosen, gain on test under the protocol, and gain on test when $E$ and $\lambda$ are fitted on test (an upper bound).}
\label{fig:sitpcl}
\end{figure}

\begin{figure}[H]
\centering
\includegraphics[width=\columnwidth]{figures/fig_gain_vs_base.pdf}
\caption{Gain of \spin{} against the strength of the base model, for all base models on both datasets. Hollow markers: not significant.}
\label{fig:gain-vs-base}
\end{figure}

\begin{figure}[H]
\centering
\includegraphics[width=\columnwidth]{figures/fig_ci_forest.pdf}
\caption{Change in Macro-F1 and \favg{} with 95\% paired-bootstrap confidence intervals (kNN few-shot base models).}
\label{fig:ci}
\end{figure}

\begin{figure}[H]
\centering
\includegraphics[width=\columnwidth]{figures/fig_ablation.pdf}
\caption{Key ablation as gains over the base model (Table~\ref{tab:ablation}).}
\label{fig:ablation}
\end{figure}

\begin{figure}[H]
\centering
\includegraphics[width=\columnwidth]{figures/fig_knowledge_type.pdf}
\caption{Relative-only, absolute-only and combined knowledge in all LLM settings (gain over the base model).}
\label{fig:knowledge-type}
\end{figure}

\begin{figure}[H]
\centering
\includegraphics[width=\columnwidth]{figures/fig_retrieval.pdf}
\caption{Relation accuracy of different ways of predicting the parent--child relation.}
\label{fig:retrieval}
\end{figure}

\begin{figure}[H]
\centering
\includegraphics[width=\columnwidth]{figures/fig_error_corr.pdf}
\caption{Parent--child error correlation of every base model (risk ratio; 1 = independent errors).}
\label{fig:errcorr}
\end{figure}

\begin{figure}[H]
\centering
\includegraphics[width=\columnwidth]{figures/fig_finetune_curve.pdf}
\caption{Qwen2.5-7B fine-tuned on 10--100\% of the ZS-CSD training split. (a) Test Macro-F1 of the model, with \spin{}, and with \spin{} given gold parent stances. (b) Accuracy of the parent--child relation implied by the model's predictions given the gold parent stance, against the retrieved relations of the knowledge base and two prompted LLMs.}
\label{fig:finetune}
\end{figure}

\section{Supplementary Tables}
\label{app:tables}

\begin{table}[H]
\centering
\caption{Selected $\lambda_r$ / $\lambda_a$ for each setting (grid $\{0, 0.25, 0.5, 0.75, 1, 1.5, 2, 3, 4, 6\}^2$, selected on validation). -- = setting not run.}
\label{tab:lambda}
\small
\begin{tabular}{@{}lcc@{}}
\toprule
Base model & C-MTCSD & ZS-CSD \\
\midrule
GPT-4o-mini kNN & 0.75 / 0.25 & 3 / 1 \\
GPT-4o kNN & 0.5 / 0.25 & 1 / 0 \\
Llama-3.1-8B kNN & 0.25 / 1 & 2 / 0.75 \\
Llama-3.1-70B kNN & 0.25 / 1 & 4 / 1.5 \\
\midrule
GPT-4o-mini CoT & 0.5 / 0.25 & 1.5 / 0.75 \\
GPT-4o-mini full-path & 0 / 0.5 & 2 / 1 \\
GPT-4o CoT & -- & 1 / 0 \\
\midrule
GLAN, seed 0 & 0.25 / 0.75 & -- \\
GLAN, seed 1 & 0.5 / 0.5 & -- \\
GLAN, seed 2 & 0.5 / 0.75 & -- \\
SITPCL, seed 1234 & -- & 2 / 1 \\
SITPCL, seed 42 & -- & 3 / 2 \\
SITPCL, seed 2026 & -- & 2 / 1.5 \\
\bottomrule
\end{tabular}
\end{table}

\begin{table}[H]
\centering
\caption{Main versus cross-validated (CV) selection of $\lambda$ (pooled Macro-F1, \%; kNN few-shot LLMs). All eight LLM settings remain significant under CV; SITPCL is not improved under either protocol.}
\label{tab:cv}
\small
\setlength{\tabcolsep}{4pt}
\begin{tabular}{@{}llccc@{}}
\toprule
 & Base model & Base & Main & CV \\
\midrule
\multirow{5}{*}{C} & GPT-4o-mini & 68.25 & 72.51 & 71.68 \\
 & GPT-4o & 74.62 & 77.78 & 77.78 \\
 & Llama-3.1-8B & 53.81 & 66.30 & 66.21 \\
 & Llama-3.1-70B & 73.41 & 74.78 & 74.38 \\
 & GLAN (3 seeds) & 68.02 & 70.05 & 70.26 \\
\midrule
\multirow{5}{*}{Z} & GPT-4o-mini & 43.71 & 46.91 & 46.91 \\
 & GPT-4o & 56.26 & 57.86 & 57.86 \\
 & Llama-3.1-8B & 29.05 & 43.52 & 42.31 \\
 & Llama-3.1-70B & 43.97 & 47.13 & 48.48 \\
 & SITPCL (seed 1234) & 45.16 & 44.84 & 43.21 \\
\bottomrule
\end{tabular}
\end{table}

\begin{table}[H]
\centering
\caption{Results on the English MT-CSD test split (pooled, \%; $n = 2{,}336$; MF1 = Macro-F1). \sig{} significant improvement over the base model ($p < 0.05$ under paired bootstrap, approximate randomization and McNemar's test).}
\label{tab:mtcsd}
\small
\setlength{\tabcolsep}{3.5pt}
\begin{tabular}{@{}lccc@{}}
\toprule
 & MF1 & \favg{} & $\lambda_r$ / $\lambda_a$ \\
\midrule
GPT-4o-mini kNN (base) & 59.91 & 57.65 & -- \\
\quad + relative only & 59.91 & 57.65 & 0 / 0 \\
\quad + absolute only & 61.08 & 58.76 & 0 / 0.75 \\
\quad + \spin{} (both) & \textbf{61.39}\sig & 58.43 & 0.75 / 0.25 \\
\midrule
\multicolumn{4}{@{}l}{\emph{Retrieval key (both kinds of knowledge)}} \\
\quad type + depth $(p,b)$ & 61.13 & 57.64 & 0.75 / 0.5 \\
\quad type only $p$ & 61.44 & 58.11 & 0.75 / 0.5 \\
\quad depth only $b$ & 59.82 & 57.44 & 0.75 / 0 \\
\quad random key & 59.91 & 57.65 & 0 / 0 \\
\midrule
\quad gold parent (upper bound) & 61.26 & 57.12 & 1 / 0 \\
\bottomrule
\end{tabular}
\end{table}

\begin{table}[H]
\centering
\caption{Gold-parent upper bound (pooled Macro-F1, \%), the numbers behind Fig.~\ref{fig:room}.}
\label{tab:room}
\small
\begin{tabular}{@{}lccc@{}}
\toprule
Setting & Base & \spin{} & gold parent \\
\midrule
C, GPT-4o-mini kNN & 68.25 & 72.51 & 71.81 \\
Z, GPT-4o-mini kNN & 43.71 & 46.91 & 58.30 \\
Z, GPT-4o-mini CoT & 46.11 & 51.49 & 59.08 \\
Z, GPT-4o CoT & 65.94 & 66.38 & 67.79 \\
Z, Qwen2.5-7B SFT & 64.55 & 64.55 & 66.09 \\
\bottomrule
\end{tabular}
\end{table}

\begin{table}[H]
\centering
\caption{Per-topic \favg{} (\%) on C-MTCSD, GPT-4o-mini kNN.}
\label{tab:topic}
\small
\begin{tabular}{@{}lccc@{}}
\toprule
Topic & Base & \spin{} & $\Delta$ \\
\midrule
iPhone 15 & 71.39 & 72.76 & +1.37 \\
Apollo Go & 67.30 & 72.83 & +5.53 \\
Pre-made Meals & 66.89 & 67.87 & +0.98 \\
Naked Resignation & 81.00 & 84.91 & +3.91 \\
Non-marriage Doctrine & 58.19 & 58.14 & $-0.05$ \\
\midrule
Average & 68.95 & 71.30 & +2.35 \\
\bottomrule
\end{tabular}
\end{table}

\begin{table}[H]
\centering
\caption{More cases (translated from Chinese and shortened). G: B$\rightarrow$S = gold: base $\rightarrow$ \spin{}. Annotations: C4 mock parent / sarcastic reversal; Z5a question parent / reversal; all others disagree / reversal. F / A / N = \favor{} / \against{} / \nonelabel{}; C = C-MTCSD, Z = ZS-CSD.}
\label{tab:more-cases}
\footnotesize
\setlength{\tabcolsep}{2.5pt}
\renewcommand{\arraystretch}{1.05}
\begin{tabular}{@{}p{0.5cm}p{2.7cm}p{2.7cm}p{1.25cm}@{}}
\toprule
ID & Parent (gold / inf.) & Reply & G: B$\rightarrow$S \\
\midrule
C4 & ``Can it guarantee it won't be hacked?'' (A / N) & ``Can you guarantee your bank card won't be hacked?'' & F: N$\rightarrow$\textbf{F}\,$\checkmark$ \\
Z4 & ``\ldots Teachers won't bother to discipline, then complain that parents make trouble when they do. Isn't that drama?'' (A / A) & ``\ldots Parents complain that teachers won't bother to discipline, then make a scene at school when they do. Isn't that drama?'' & F: A$\rightarrow$\textbf{F}\,$\checkmark$ \\
\midrule
C5a & Root post: a nurse quits a top hospital without another job. (F / N) & ``Yet thousands of graduates queue up to get in every year.'' & N: N$\rightarrow$\textbf{A}\,$\times$ \\
C5b & C5a (N / A) & ``Should the hospitals just close, then? Where would you see a doctor?'' & N: N$\rightarrow$\textbf{A}\,$\times$ \\
Z5a & ``This is acquired, caused by family environment and upbringing.'' (N / A) & ``So upbringing can turn a good gene into a criminal, but cannot fix a suspected bad gene?'' & A: A$\rightarrow$\textbf{F}\,$\times$ \\
Z5b & Z5a (A / F) & ``What kind of logic is that? A classmate of mine \ldots{} twenty years later he is a normal adult.'' & N: N$\rightarrow$\textbf{A}\,$\times$ \\
\bottomrule
\end{tabular}
\end{table}

\begin{table}[H]
\centering
\caption{Fine-tuned Qwen2.5-7B on ZS-CSD (test Macro-F1, \%). The knowledge base uses the same training subset. \sig{} significant decrease. Data: share of the training split. Gold par.: \spin{} with gold parent stances. Rel.\ acc.: accuracy of the parent--child relation implied by the predictions given the gold parent stance (\%); for reference, 55.0 for the relations retrieved from the knowledge base, 41.7 for GPT-4o-mini kNN and 34.8 for Llama-3.1-8B kNN.}
\label{tab:finetune}
\small
\setlength{\tabcolsep}{3pt}
\begin{tabular}{@{}lccccc@{}}
\toprule
Data & Base & +\spin{} & $\Delta$ & Gold par. & Rel.\ acc. \\
\midrule
10\% & 53.35 & 52.07 & $-1.27$ & 57.92 & 51.3 \\
25\% & 59.17 & 58.12 & $-1.06$ & 62.90 & 56.5 \\
50\% & 61.27 & 60.64 & $-0.63$\sig & 65.43 & 58.0 \\
100\% & 64.55 & 64.55 & 0.00 & 66.09 & 61.3 \\
\bottomrule
\end{tabular}
\end{table}

\begin{table}[H]
\centering
\caption{\spin{} versus combinations learned on the validation split from the same inputs (pooled test Macro-F1, \%; kNN few-shot base models). $^{+}$/$^{-}$: significantly better/worse than \spin{} (paired bootstrap and approximate randomization, $p<0.05$).}
\label{tab:learned}
\small
\setlength{\tabcolsep}{3.5pt}
\begin{tabular}{@{}llcccc@{}}
\toprule
 & Base model & Base & \spin{} & LR-log & GBDT \\
\midrule
C & GPT-4o-mini & 68.25 & 72.51 & 72.61 & 72.37 \\
C & Llama-3.1-8B & 53.81 & 66.30 & 66.73 & \textbf{67.54}$^{+}$ \\
Z & GPT-4o-mini & 43.71 & \textbf{46.91} & 45.67 & 44.59$^{-}$ \\
Z & Llama-3.1-8B & 29.05 & \textbf{43.52} & 40.04$^{-}$ & 43.42 \\
\bottomrule
\end{tabular}
\end{table}

\begin{table}[H]
\centering
\caption{Cost of the semantic annotation over all splits (GPT-4o-mini list prices; input tokens counted with the GPT-4o tokenizer, output estimated at 250 tokens per call). Root posts are annotated by rule without a call.}
\label{tab:cost}
\small
\setlength{\tabcolsep}{4pt}
\begin{tabular}{@{}lrrrr@{}}
\toprule
Dataset & Nodes & LLM calls & Input tok./call & US\$ \\
\midrule
C-MTCSD & 24,263 & 23,811 & 1,239 & 8.0 \\
ZS-CSD & 17,008 & 16,709 & 1,205 & 5.5 \\
MT-CSD & 15,287 & 15,287 & 1,154 & 4.9 \\
\bottomrule
\end{tabular}
\end{table}

\section{Annotation Prompt}
\label{app:prompt}

The annotation prompt, translated from Chinese. Slots in braces are filled per node. In the causal mode used in this paper, the downstream-reply and post-hoc slots are replaced by the fixed strings shown, so the fields that refer to child replies are always \texttt{no\_child} / \texttt{false}. \spin{} uses only \texttt{llm\_reply\_relation}, \texttt{llm\_semantic\_propagation} and \texttt{llm\_target\_relation}; all other fields, including the stance guess, are ignored.

\begin{lstlisting}[style=prompt]
=== SYSTEM PROMPT ===
You are an expert in stance detection and semantic propagation analysis for social-media conversations. Your task is not simply to judge sentiment, but to analyse the semantic function of the current comment in the conversation chain, its reply relation, its target orientation, and how stance may propagate. You must strictly distinguish the main_target, the parent comment, the author being replied to, and topic shifts. Output only JSON, nothing else.

=== USER PROMPT TEMPLATE ===
Analyse the following social-media conversation sample.

Main target main_target: {main_target}
Topic: {topic}
Target anchor target_anchor: {target_anchor}
Upstream conversation window upward_window: {upward_window_texts}
Current comment current_text: {current_text}
Downstream reply feedback downward_feedback: (causal mode: replies are not used)
Full causal_context_text: {causal_context_text}
Full posthoc_context_text: (causal mode: not provided)

Please complete the following analysis:
1. Judge the reply relation between the current comment and its parent comment.
2. Judge whether meaning propagates from the parent to the current comment: inheritance, reversal, amplification, softening, sarcastic reversal, target shift, or unclear.
3. Judge how the object directly discussed by the current comment relates to main_target.
4. If the current comment is short or ambiguous, judge whether child replies help clarify its meaning.
5. Infer the likely stance of the current comment toward main_target (this is AI-silver, not human gold).
6. List the linguistic structure patterns.
7. Give a confidence (0-1) and a one-sentence reason (at most 80 characters).

Important requirements:
- Do not simply treat "irrelevant", "unclear target" or "insufficient context" as none.
- If the current comment is a short text such as "indeed", "heh", "666", "lol" or "classic", you must use the parent comment to judge its pragmatic function.
- If a child reply says things like "are you praising or mocking?", "obviously sarcastic" or "you misunderstood", it may be used to infer the true meaning of the current comment; mark llm_use_child_feedback as true.
- If the current comment only attacks the author being replied to and does not discuss main_target, set llm_target_relation to reply_author_only.
- If the current comment discusses an attribute, impact, consequence or implementation of main_target, use attribute_of_target or consequence_of_target rather than topic_shift.
- Output valid JSON only, nothing else.

Output strictly in the following JSON schema (all fields must be present):
{
 "llm_reply_relation": "agree_with_parent|disagree_with_parent|mock_parent|sarcastic_agreement|clarify_parent|supplement_parent|question_parent|answer_parent|topic_shift|direct_comment_to_target|unknown",
 "llm_semantic_propagation": "stance_inheritance|stance_reversal|sarcastic_reversal|stance_amplification|stance_softening|direct_stance_expression|target_shift|unclear_propagation|no_parent_context",
 "llm_stance_mapping_to_parent": "same_as_parent|opposite_to_parent|partial_or_conditional|independent_from_parent|unclear|not_applicable",
 "llm_target_relation": "same_target|attribute_of_target|consequence_of_target|reason_for_target|opposition_to_related_target|support_to_related_target|reply_author_only|topic_shift|irrelevant|unclear",
 "llm_child_feedback_signal": "child_confirms_current_meaning|child_questions_current_meaning|child_corrects_misreading|child_indicates_sarcasm|child_indicates_disagreement|child_conflict|no_useful_child_feedback|no_child",
 "llm_candidate_patterns": ["agreement_with_parent|disagreement_with_parent|mocking_disagreement|...|unknown_pattern"],
 "llm_direct_target": "main_target|parent_comment|reply_author|target_attribute|target_consequence|related_entity|unclear",
 "llm_inferred_stance_to_main_target": "favor|against|none|mixed_or_conditional|unclear|irrelevant",
 "llm_confidence": 0.0,
 "llm_reason_short": "a reason of at most 80 characters",
 "llm_use_child_feedback": false,
 "llm_context_setting": "causal|posthoc|both"
}
\end{lstlisting}